\documentclass[12pt,a4paper]{article}
\usepackage[a4paper,margin=1.8cm]{geometry}
\usepackage[T2A]{fontenc}
\usepackage[utf8]{inputenc}
\usepackage[russian]{babel}
\usepackage{amsmath,amssymb,mathtools}
\usepackage{array,booktabs,longtable}
\usepackage{xcolor}
\usepackage{hyperref}
\usepackage{tikz}
\hypersetup{colorlinks=true,linkcolor=blue!55!black,urlcolor=blue!55!black}
\setlength{\parindent}{0pt}
\setlength{\parskip}{5pt}
\renewcommand{\arraystretch}{1.18}
\newcommand{\err}{\textcolor{red}{Ошибка:}}
\newcommand{\idea}{\textcolor{blue!60!black}{Ключевая идея:}}
\newcommand{\result}{\textcolor{teal!60!black}{Итог:}}
\newcommand{\tasktitle}[1]{\phantomsection\label{task:#1}\section*{Задание №#1}}

\begin{document}

\begin{titlepage}
\centering
\vspace*{0.22\textheight}
{\LARGE\bfseries Ревью домашней работы\par}
\vspace{1.6cm}
{\large Автор: Лёвин Артём Александрович\par}
\vspace{0.5cm}
{\large 29 сентября 2026 г.\par}
\vfill
\begin{tikzpicture}[scale=1]
\draw[blue!45!black, line width=0.7pt] (0,0) .. controls (2,0.65) and (4,-0.65) .. (6,0);
\end{tikzpicture}
\end{titlepage}

\newpage
\section*{Сводная таблица}
\small
\begin{longtable}{>{\centering\arraybackslash}p{0.8cm} p{3.0cm} p{5.0cm} p{2.6cm} p{2.9cm}}
\toprule
№ & Тема & Суть ошибки & Ваш ответ & Правильный ответ \\
\midrule
\endfirsthead
\toprule
№ & Тема & Суть ошибки & Ваш ответ & Правильный ответ \\
\midrule
\endhead
\hyperref[task:5]{5} & Теорема Виета & Вместо использования сумм и произведения корней начато прямое решение квадратного уравнения; требуемая сумма квадратов корней не найдена. & Требуемый ответ не указан & $54$ \\
\addlinespace
\hyperref[task:9]{9} & Показательное уравнение & Проверен только один подходящий корень. Второй корень уравнения потерян. & $2$ & $0;\ 2$ \\
\addlinespace
\hyperref[task:10]{10} & Логарифмы & Не учтено условие существования логарифмов, поэтому в ответ оставлен посторонний корень; также в промежуточной строке сумма логарифмов неверно записана как сумма аргументов. & $5;\ -1$ & $5$ \\
\addlinespace
\hyperref[task:13]{13} & Область определения & Решение в материалах не представлено. & Не указан & $(-\infty;2)\cup(2;3]$ \\
\addlinespace
\hyperref[task:14]{14} & Квадратичная функция & Решение в материалах не представлено. & Не указан & $-4$ \\
\addlinespace
\hyperref[task:15]{15} & Арифметическая прогрессия & Решение в материалах не представлено. & Не указан & $210$ \\
\addlinespace
\hyperref[task:16]{16} & Производная и монотонность & Решение в материалах не представлено. & Не указан & $(-\infty;-1)$ и $(3;+\infty)$ \\
\addlinespace
\hyperref[task:23]{23} & Сложные проценты & Решение в материалах не представлено. & Не указан & $10\%$ \\
\addlinespace
\hyperref[task:24]{24} & Квадратное уравнение с параметром & Решение в материалах не представлено. & Не указан & $a=-1$ или $a=2$ \\
\bottomrule
\end{longtable}
\normalsize

\newpage
\vspace*{0.32\textheight}
\begin{center}
{\LARGE\bfseries Разбор ошибок}
\end{center}
\vfill
\newpage

\tasktitle{5}
\textbf{Текст задания.} Числа $x_1$ и $x_2$ --- корни уравнения
\[
x^2-8x+5=0.
\]
Не решая уравнение напрямую, найдите $x_1^2+x_2^2$.

\textbf{Ваше решение.} Выписан дискриминант
\[
D=(-8)^2-4\cdot1\cdot5=44,
\]
после чего начато вычисление самих корней по формуле
\[
x_{1,2}=\frac{8\pm\sqrt{44}}{2}.
\]

\textbf{Ваш ответ.} Требуемое значение $x_1^2+x_2^2$ не указано.

\err{} выбран путь прямого решения квадратного уравнения, хотя условие специально просит не находить корни. Кроме того, решение остановлено до вычисления требуемого выражения.

\textbf{Где именно возникает ошибка.} Последний полностью корректный шаг --- вычисление дискриминанта $D=44$. Первый неполный шаг --- переход к формуле корней вместо использования связей между корнями и коэффициентами.

\textbf{Почему это неверно.} Для приведённого квадратного уравнения $x^2+bx+c=0$ теорема Виета сразу даёт сумму и произведение корней. В данном случае
\[
x_1+x_2=8,\qquad x_1x_2=5.
\]
Искомое выражение зависит именно от этих двух величин, поэтому сами корни находить не требуется. Прямое вычисление корней не только нарушает требование условия, но и создаёт лишнюю работу с радикалами.

\idea{} квадрат суммы корней связывает искомую сумму квадратов с произведением корней:
\[
(x_1+x_2)^2=x_1^2+2x_1x_2+x_2^2.
\]
Отсюда
\[
x_1^2+x_2^2=(x_1+x_2)^2-2x_1x_2.
\]

\textbf{Исправление от последнего правильного шага.} Вместо дальнейшего вычисления корней используем коэффициенты исходного уравнения:
\[
x_1+x_2=8,\qquad x_1x_2=5.
\]
Тогда
\[
x_1^2+x_2^2=8^2-2\cdot5=64-10=54.
\]

\textbf{Полное правильное решение.} По теореме Виета для уравнения $x^2-8x+5=0$ имеем
\[
x_1+x_2=8,\qquad x_1x_2=5.
\]
Используем тождество
\[
x_1^2+x_2^2=(x_1+x_2)^2-2x_1x_2.
\]
Подставляем известные значения:
\[
x_1^2+x_2^2=8^2-2\cdot5=54.
\]

\textbf{Проверка.} Формула получена непосредственным раскрытием квадрата суммы, поэтому дополнительных ограничений нет. Числа $8$ и $5$ действительно являются соответственно суммой и произведением корней данного уравнения.

\result{} $x_1^2+x_2^2=54$.

\subsection*{Задача на закрепление}
Числа $u_1$ и $u_2$ --- корни уравнения $u^2-10u+7=0$. Не решая уравнение напрямую, найдите $u_1^2+u_2^2$.

\newpage
\tasktitle{9}
\textbf{Текст задания.} Решите уравнение
\[
4^x-5\cdot2^x+4=0.
\]

\textbf{Ваше решение.} Подстановкой проверено значение $x=2$:
\[
16-20+4=0,
\]
после чего записан ответ $2$.

\textbf{Ваш ответ.} $2$.

\err{} найден только один корень, а уравнение имеет два действительных решения.

\textbf{Где именно возникает ошибка.} Проверка $x=2$ выполнена верно. Первый неверный шаг --- считать, что после нахождения одного подходящего значения уравнение решено полностью.

\textbf{Почему это неверно.} Выражения $4^x$ и $2^x$ связаны равенством $4^x=(2^x)^2$. Поэтому исходное уравнение фактически является квадратным относительно величины $2^x$. Квадратное уравнение может иметь два допустимых корня, и проверка одного значения не доказывает отсутствие другого.

\idea{} ввести замену $t=2^x$. Поскольку $2^x>0$ при любом действительном $x$, получаем обычное квадратное уравнение с дополнительным условием $t>0$.

\textbf{Исправление от последнего правильного шага.} Значение $x=2$ сохраняем как найденный корень, но продолжаем решать уравнение системно. Положим
\[
t=2^x,\qquad t>0.
\]
Тогда
\[
t^2-5t+4=0.
\]
Разложим на множители:
\[
(t-1)(t-4)=0.
\]
Получаем
\[
t=1\quad\text{или}\quad t=4.
\]
Возвращаемся к $x$:
\[
2^x=1\Rightarrow x=0,
\]
\[
2^x=4\Rightarrow x=2.
\]

\textbf{Полное правильное решение.} Используем $4^x=(2^x)^2$:
\[
(2^x)^2-5\cdot2^x+4=0.
\]
Пусть $t=2^x$, тогда $t>0$. Получаем
\[
t^2-5t+4=0,
\]
откуда
\[
(t-1)(t-4)=0.
\]
Следовательно,
\[
t=1\quad\text{или}\quad t=4.
\]
Значит,
\[
2^x=1\Rightarrow x=0,
\qquad
2^x=4\Rightarrow x=2.
\]

\textbf{Проверка.} При $x=0$:
\[
1-5+4=0.
\]
При $x=2$:
\[
16-20+4=0.
\]
Оба значения подходят.

\result{} $x=0$ или $x=2$.

\subsection*{Задача на закрепление}
Решите уравнение $9^x-10\cdot3^x+9=0$.

\newpage
\tasktitle{10}
\textbf{Текст задания.} Решите уравнение
\[
\log_2(x-1)+\log_2(x-3)=3.
\]

\textbf{Ваше решение.} Правая часть записана как $\log_2 8$. Далее получено квадратное уравнение
\[
x^2-4x-5=0
\]
и корни $5$ и $-1$. В одной промежуточной строке между выражениями $x-1$ и $x-3$ записан знак сложения, хотя последующее раскрытие соответствует их произведению.

\textbf{Ваш ответ.} В записи оставлены значения $5$ и $-1$.

\err{} не учтена область допустимых значений логарифмов, поэтому значение $-1$ не отброшено. Кроме того, сумма логарифмов должна превращаться в логарифм произведения аргументов, а не их суммы.

\textbf{Где именно возникает ошибка.} Корректно замечено, что $3=\log_2 8$. Первый обязательный шаг, который пропущен, --- запись условий
\[
x-1>0,\qquad x-3>0.
\]
Они дают $x>3$. После решения квадратного уравнения это условие нужно проверить для каждого найденного корня.

\textbf{Почему это неверно.} Логарифм определён только при положительном аргументе. Поэтому исходное уравнение вообще не имеет смысла при $x\leq3$. Формула сложения логарифмов имеет вид
\[
\log_2 A+\log_2 B=\log_2(AB)
\]
при $A>0$ и $B>0$. Следовательно, аргументы нужно перемножать.

\idea{} сначала выписать область допустимых значений, затем объединить логарифмы и только после этого решать алгебраическое уравнение.

\textbf{Исправление от последнего правильного шага.} Записываем условия существования логарифмов:
\[
x>1,\qquad x>3,
\]
то есть
\[
x>3.
\]
Далее
\[
\log_2\bigl((x-1)(x-3)\bigr)=\log_2 8.
\]
Так как основания одинаковы,
\[
(x-1)(x-3)=8.
\]
Раскрываем скобки:
\[
x^2-4x+3=8,
\]
\[
x^2-4x-5=0.
\]
Корни:
\[
x=5,\qquad x=-1.
\]
С учётом условия $x>3$ остаётся только $x=5$.

\textbf{Полное правильное решение.} Область допустимых значений:
\[
x-1>0,\qquad x-3>0,
\]
следовательно,
\[
x>3.
\]
Объединяем логарифмы:
\[
\log_2\bigl((x-1)(x-3)\bigr)=3=\log_2 8.
\]
Отсюда
\[
(x-1)(x-3)=8.
\]
Получаем
\[
x^2-4x-5=0,
\]
\[
(x-5)(x+1)=0.
\]
Поэтому возможны $x=5$ и $x=-1$, но $x=-1$ не принадлежит области допустимых значений.

\textbf{Проверка.} При $x=5$:
\[
\log_2 4+\log_2 2=2+1=3.
\]
Значение $x=-1$ подставлять в исходное выражение нельзя, поскольку аргументы логарифмов неположительны.

\result{} $x=5$.

\subsection*{Задача на закрепление}
Решите уравнение $\log_3(x-2)+\log_3(x-4)=2$.

\newpage
\tasktitle{13}
\textbf{Текст задания.} Найдите область определения функции
\[
f(x)=\frac{\sqrt{6-2x}}{x-2}.
\]

\textbf{Ваше решение.} В представленных материалах решение этого задания отсутствует.

\textbf{Ваш ответ.} Не указан.

\err{} не выполнены обязательные условия существования квадратного корня и дроби.

\textbf{Где именно возникает ошибка.} Последнего корректного шага в записи нет, поскольку решение не начато. Первый необходимый шаг --- одновременно потребовать неотрицательность подкоренного выражения и ненулевой знаменатель.

\textbf{Почему это неверно.} Для квадратного корня необходимо
\[
6-2x\geq0.
\]
Для дроби дополнительно требуется
\[
x-2\ne0.
\]
Область определения должна удовлетворять обоим условиям одновременно.

\idea{} область определения составляется пересечением всех ограничений, возникающих из отдельных частей формулы.

\textbf{Исправление.} Из первого условия
\[
6-2x\geq0
\]
получаем
\[
-2x\geq-6,
\]
а при делении на отрицательное число знак неравенства меняется:
\[
x\leq3.
\]
Из второго условия
\[
x\ne2.
\]
Значит, из всех чисел, не превосходящих $3$, нужно исключить $2$.

\textbf{Полное правильное решение.} Требования к выражению:
\[
\begin{cases}
6-2x\geq0,\\
x-2\ne0.
\end{cases}
\]
Первое условие даёт $x\leq3$, второе исключает $x=2$. Поэтому
\[
D(f)=(-\infty;2)\cup(2;3].
\]

\textbf{Проверка.} При $x=3$ подкоренное выражение равно нулю, знаменатель равен единице, поэтому точка $3$ входит в область определения. При $x=2$ знаменатель обращается в ноль, поэтому точка $2$ исключается.

\result{} $D(f)=(-\infty;2)\cup(2;3]$.

\subsection*{Задача на закрепление}
Найдите область определения функции $g(x)=\dfrac{\sqrt{8-2x}}{x+1}$.

\newpage
\tasktitle{14}
\textbf{Текст задания.} Найдите наименьшее значение функции
\[
y=x^2-6x+5.
\]

\textbf{Ваше решение.} В представленных материалах решение этого задания отсутствует.

\textbf{Ваш ответ.} Не указан.

\err{} не найдено минимальное значение квадратного трёхчлена.

\textbf{Где именно возникает ошибка.} Решение не начато. Первый необходимый шаг --- привести квадратный трёхчлен к виду квадрата двучлена плюс постоянное число либо найти вершину параболы.

\textbf{Почему это важно.} Коэффициент при $x^2$ положителен, поэтому ветви параболы направлены вверх. Следовательно, значение функции в вершине является наименьшим.

\idea{} выделить полный квадрат:
\[
x^2-6x=(x-3)^2-9.
\]

\textbf{Исправление.} Преобразуем:
\[
y=x^2-6x+5=(x-3)^2-9+5=(x-3)^2-4.
\]
Квадрат любого действительного числа неотрицателен:
\[
(x-3)^2\geq0.
\]
Поэтому
\[
y\geq-4.
\]
Равенство достигается при $x=3$.

\textbf{Полное правильное решение.} Выделяем полный квадрат:
\[
y=(x-3)^2-4.
\]
Так как $(x-3)^2\geq0$, наименьшее возможное значение выражения получается при $(x-3)^2=0$, то есть при $x=3$. Тогда
\[
y_{\min}=-4.
\]

\textbf{Проверка.} При $x=3$ исходная функция равна
\[
3^2-6\cdot3+5=9-18+5=-4.
\]

\result{} наименьшее значение функции равно $-4$.

\subsection*{Задача на закрепление}
Найдите наименьшее значение функции $y=x^2-8x+11$.

\newpage
\tasktitle{15}
\textbf{Текст задания.} В арифметической прогрессии $a_4=10$, $a_9=25$. Найдите сумму первых двенадцати членов $S_{12}$.

\textbf{Ваше решение.} В представленных материалах решение этого задания отсутствует.

\textbf{Ваш ответ.} Не указан.

\err{} не использованы данные о двух членах прогрессии для нахождения разности и первого члена.

\textbf{Где именно возникает ошибка.} Решение не начато. Первый необходимый шаг --- записать формулу $a_n=a_1+(n-1)d$ для двух известных номеров.

\textbf{Почему это работает.} Все члены арифметической прогрессии определяются первым членом $a_1$ и постоянной разностью $d$. Два независимых условия позволяют найти обе неизвестные величины.

\idea{} составить систему
\[
a_1+3d=10,
\qquad
a_1+8d=25.
\]

\textbf{Исправление.} Вычитаем первое равенство из второго:
\[
5d=15,
\]
поэтому
\[
d=3.
\]
Тогда
\[
a_1+3\cdot3=10,
\]
откуда
\[
a_1=1.
\]
Сумма первых двенадцати членов:
\[
S_{12}=\frac{12}{2}\bigl(2a_1+11d\bigr).
\]
Подставляем:
\[
S_{12}=6(2+33)=6\cdot35=210.
\]

\textbf{Полное правильное решение.} По формуле общего члена
\[
a_4=a_1+3d=10,
\]
\[
a_9=a_1+8d=25.
\]
Вычитая, получаем $5d=15$, значит $d=3$. Тогда $a_1=1$. Используем формулу суммы:
\[
S_{12}=\frac{12}{2}\bigl(2\cdot1+11\cdot3\bigr)=6\cdot35=210.
\]

\textbf{Проверка.} При $a_1=1$ и $d=3$ имеем
\[
a_4=1+9=10,
\qquad
a_9=1+24=25,
\]
то есть исходные данные выполняются.

\result{} $S_{12}=210$.

\subsection*{Задача на закрепление}
В арифметической прогрессии $a_3=8$, $a_7=20$. Найдите сумму первых десяти членов.

\newpage
\tasktitle{16}
\textbf{Текст задания.} Для функции
\[
f(x)=x^3-3x^2-9x+7
\]
найдите промежутки возрастания.

\textbf{Ваше решение.} В представленных материалах решение этого задания отсутствует.

\textbf{Ваш ответ.} Не указан.

\err{} не выполнено исследование знака производной.

\textbf{Где именно возникает ошибка.} Решение не начато. Первый необходимый шаг --- найти $f'(x)$ и определить, где производная положительна.

\textbf{Почему это работает.} На промежутках, где $f'(x)>0$, функция возрастает. Поэтому задача сводится к решению неравенства для производной.

\idea{} производная имеет вид
\[
f'(x)=3x^2-6x-9,
\]
и удобно вынести общий множитель и разложить квадратный трёхчлен на множители.

\textbf{Исправление.} Находим производную:
\[
f'(x)=3x^2-6x-9=3(x^2-2x-3)=3(x-3)(x+1).
\]
Критические точки:
\[
x=-1,\qquad x=3.
\]
Знак произведения положителен вне промежутка между корнями, поэтому
\[
f'(x)>0
\]
при
\[
x<-1\quad\text{или}\quad x>3.
\]

\textbf{Полное правильное решение.} Вычисляем производную:
\[
f'(x)=3x^2-6x-9=3(x-3)(x+1).
\]
Решаем неравенство
\[
3(x-3)(x+1)>0.
\]
Поскольку множитель $3$ положителен, знак определяется произведением $(x-3)(x+1)$. Оно положительно, когда оба множителя одного знака. Получаем
\[
x<-1\quad\text{или}\quad x>3.
\]
Следовательно, функция возрастает на промежутках
\[
(-\infty;-1)\quad\text{и}\quad(3;+\infty).
\]

\textbf{Проверка.} Например, при $x=-2$ производная положительна, при $x=0$ отрицательна, при $x=4$ снова положительна. Это подтверждает найденную схему знаков.

\result{} промежутки возрастания: $(-\infty;-1)$ и $(3;+\infty)$.

\subsection*{Задача на закрепление}
Для функции $g(x)=x^3-6x^2+9x+4$ найдите промежутки возрастания.

\newpage
\tasktitle{23}
\textbf{Текст задания.} Вклад увеличился со $100\,000$ рублей до $121\,000$ рублей за два года. Проценты начислялись один раз в год по одной и той же ставке и каждый раз прибавлялись к сумме вклада. Найдите годовую процентную ставку.

\textbf{Ваше решение.} В представленных материалах решение этого задания отсутствует.

\textbf{Ваш ответ.} Не указан.

\err{} не составлена модель двух последовательных начислений сложных процентов.

\textbf{Где именно возникает ошибка.} Решение не начато. Первый необходимый шаг --- обозначить годовой коэффициент увеличения и учесть, что он применяется два раза подряд.

\textbf{Почему это важно.} После первого года проценты начисляются уже не на исходные $100\,000$ рублей, а на увеличенную сумму. Поэтому общий рост за два года равен квадрату годового коэффициента, а не удвоенному проценту.

\idea{} если годовая ставка равна $r$ в долях единицы, то после каждого года сумма умножается на $1+r$.

\textbf{Исправление.} Составляем уравнение:
\[
100000(1+r)^2=121000.
\]
Делим на $100000$:
\[
(1+r)^2=1.21.
\]
Так как $1+r$ --- положительный коэффициент роста,
\[
1+r=1.1.
\]
Следовательно,
\[
r=0.1.
\]
В процентах это $10\%$.

\textbf{Полное правильное решение.} Пусть годовая ставка равна $r$. Тогда после первого года вклад станет
\[
100000(1+r),
\]
а после второго
\[
100000(1+r)^2.
\]
По условию эта сумма равна $121000$:
\[
100000(1+r)^2=121000.
\]
Отсюда
\[
(1+r)^2=1.21,
\]
\[
1+r=1.1,
\]
\[
r=0.1=10\%.
\]

\textbf{Проверка.} После первого года при ставке $10\%$ получается $110\,000$ рублей, после второго --- $121\,000$ рублей. Это совпадает с условием.

\result{} годовая процентная ставка равна $10\%$.

\subsection*{Задача на закрепление}
Вклад увеличился с $200\,000$ рублей до $242\,000$ рублей за два года при одинаковой ежегодной ставке с капитализацией. Найдите годовую процентную ставку.

\newpage
\tasktitle{24}
\textbf{Текст задания.} Найдите все значения параметра $a$, при которых уравнение
\[
x^2-2ax+a+2=0
\]
имеет ровно один действительный корень.

\textbf{Ваше решение.} В представленных материалах решение этого задания отсутствует.

\textbf{Ваш ответ.} Не указан.

\err{} не использовано условие единственного действительного корня квадратного уравнения.

\textbf{Где именно возникает ошибка.} Решение не начато. Первый необходимый шаг --- заметить, что коэффициент при $x^2$ равен единице при любом $a$, поэтому уравнение всегда остаётся квадратным, а один действительный корень возникает ровно при нулевом дискриминанте.

\textbf{Почему это верно.} Для квадратного уравнения $Ax^2+Bx+C=0$ при $A\ne0$ число действительных корней определяется дискриминантом. Ровно один действительный корень получается при $D=0$.

\idea{} вычислить дискриминант как выражение от параметра $a$ и решить уравнение $D=0$.

\textbf{Исправление.} Здесь
\[
A=1,\qquad B=-2a,\qquad C=a+2.
\]
Дискриминант:
\[
D=(-2a)^2-4\cdot1\cdot(a+2).
\]
Упрощаем:
\[
D=4a^2-4a-8=4(a^2-a-2).
\]
Для единственного корня требуется
\[
D=0.
\]
Следовательно,
\[
a^2-a-2=0.
\]
Разложим на множители:
\[
(a-2)(a+1)=0.
\]
Поэтому
\[
a=2\quad\text{или}\quad a=-1.
\]

\textbf{Полное правильное решение.} Так как коэффициент при $x^2$ всегда равен $1$, уравнение является квадратным для любого действительного $a$. Оно имеет ровно один действительный корень тогда и только тогда, когда
\[
D=0.
\]
Вычисляем:
\[
D=(-2a)^2-4(a+2)=4a^2-4a-8.
\]
Приравниваем к нулю и делим на $4$:
\[
a^2-a-2=0.
\]
Отсюда
\[
(a-2)(a+1)=0,
\]
поэтому
\[
a=2\quad\text{или}\quad a=-1.
\]

\textbf{Проверка.} При $a=2$ получаем
\[
x^2-4x+4=0,
\]
то есть $(x-2)^2=0$ --- один действительный корень. При $a=-1$ получаем
\[
x^2+2x+1=0,
\]
то есть $(x+1)^2=0$ --- также один действительный корень.

\result{} $a=-1$ или $a=2$.

\subsection*{Задача на закрепление}
Найдите все значения параметра $b$, при которых уравнение $x^2-2bx+2b-3=0$ имеет ровно один действительный корень.

\vfill
\begin{center}
\small Лёвин Артём Александрович эксклюзивно для Даниила Кичука
\end{center}

\end{document}
